\exer{Géométrie}

\setcounter{exo}{0}
Soit le schéma cinématique précédent.
\begin{center}
\includegraphics[width=\linewidth]{057_01}
\end{center}

\question{Exprimer les coordonnées du point P en fonction des angles phi et theta.}
\ifprof
\begin{corrige}

\end{corrige}
\else
\fi

\question{Exprimer les angles phi et theta en fonction des autres paramètres.}
\ifprof
\begin{corrige}

\end{corrige}
\else
\fi

\question{Donner les valeurs angulaires de phi et theta pour que le point P suive une ligne droite du point de coordonnées (L, -h) au point de coordonnées (L, h). Cette question pourra nécessiter l'utilisation de Python pour réaliser les tracés. On pourra prendre a = b = 1, L = 1 et h = 1.}
\ifprof
\begin{corrige}
\end{corrige}
\else
\fi

% BEGIN AUTO CORRIGE INCLUDE
\InclureCorrige{corrige.tex}
% END AUTO CORRIGE INCLUDE
